% notation.tex -- all notation of the document is declared here.
% Every macro is semantic: change the rendering here, not in the text.

% ---------------------------------------------------------------------
% Packages
% ---------------------------------------------------------------------
\usepackage{amsmath,amssymb,amsthm,mathtools}
\usepackage{etoolbox}
\usepackage{array}
\usepackage{tikz-cd}
\newcolumntype{L}[1]{>{\raggedright\arraybackslash}p{#1}}
\usepackage[colorinlistoftodos]{todonotes}
\usepackage[hidelinks]{hyperref}
\usepackage[capitalise,nameinlink]{cleveref}

% ---------------------------------------------------------------------
% Theorem environments
% ---------------------------------------------------------------------
\theoremstyle{plain}
\newtheorem{theorem}{Theorem}[section]
\newtheorem{proposition}[theorem]{Proposition}
\newtheorem{lemma}[theorem]{Lemma}
\newtheorem{corollary}[theorem]{Corollary}
\theoremstyle{definition}
\newtheorem{definition}[theorem]{Definition}
\newtheorem{setup}[theorem]{Setup}
\newtheorem{convention}[theorem]{Convention}
\newtheorem{example}[theorem]{Example}
\theoremstyle{remark}
\newtheorem{remark}[theorem]{Remark}

\crefname{setup}{Setup}{Setups}
\crefname{convention}{Convention}{Conventions}

% Stacks project style slogan, used as the first line of a statement.
\newcommand{\slogan}[1]{\textit{Slogan: #1}\par\smallskip\noindent}

% ---------------------------------------------------------------------
% Number systems, rings, spaces
% ---------------------------------------------------------------------
\newcommand{\NN}{\mathbb{N}}
\newcommand{\ZZ}{\mathbb{Z}}
\newcommand{\QQ}{\mathbb{Q}}
\newcommand{\RR}{\mathbb{R}}
\newcommand{\CC}{\mathbb{C}}
\newcommand{\FF}[1]{\mathbb{F}_{#1}}
\newcommand{\pAdicIntegers}[1]{\ZZ_{#1}}
\newcommand{\pAdicNumbers}[1]{\QQ_{#1}}
\newcommand{\ProfiniteIntegers}{\widehat{\ZZ}}
\newcommand{\Adeles}[1]{\mathbb{A}_{#1}}
\newcommand{\FiniteAdeles}[1]{\mathbb{A}_{#1,f}}
\newcommand{\PositiveRationals}{\QQ_{>0}}
\newcommand{\unit}[1]{#1^{\times}}
\newcommand{\ResidueRing}[1]{\ZZ/#1\ZZ}
\newcommand{\AffineSpace}[2]{\mathbb{A}^{#2}_{#1}}
\newcommand{\ProjectiveSpace}[2]{\mathbb{P}^{#2}_{#1}}
\newcommand{\ProjectiveLine}[1]{\mathbb{P}^{1}(#1)}
\newcommand{\UpperHalfPlane}{\mathfrak{H}}
\newcommand{\SlitPlane}{\CC\setminus\RR}
\newcommand{\ExtendedUpperHalfPlane}{\mathfrak{H}^{*}}
\newcommand{\RootOfUnity}[1]{\zeta_{#1}}
\newcommand{\CyclotomicField}[1]{\QQ(\RootOfUnity{#1})}
\newcommand{\MaxAbelian}{\QQ^{\mathrm{ab}}}
\newcommand{\StdBasis}[1]{e_{#1}}
\newcommand{\StdModule}[2]{#1^{#2}}

% ---------------------------------------------------------------------
% Matrix groups and matrix operations
% ---------------------------------------------------------------------
% An empty second argument gives the group scheme name, e.g. \GL{2}{}.
\newcommand{\optring}[1]{\ifstrempty{#1}{}{(#1)}}
\newcommand{\GL}[2]{\mathrm{GL}_{#1}\optring{#2}}
\newcommand{\SL}[2]{\mathrm{SL}_{#1}\optring{#2}}
\newcommand{\PGL}[2]{\mathrm{PGL}_{#1}\optring{#2}}
\newcommand{\PSL}[2]{\mathrm{PSL}_{#1}\optring{#2}}
\newcommand{\GLplus}[2]{\mathrm{GL}_{#1}^{+}(#2)}
\newcommand{\PGLplus}[2]{\mathrm{PGL}_{#1}^{+}(#2)}
\newcommand{\MatRing}[2]{\mathrm{M}_{#1}(#2)}
\newcommand{\Mat}[4]{\begin{pmatrix} #1 & #2 \\ #3 & #4 \end{pmatrix}}
\newcommand{\smallMat}[4]{\left(\begin{smallmatrix} #1 & #2 \\ #3 & #4 \end{smallmatrix}\right)}
\newcommand{\ColVec}[2]{\begin{pmatrix} #1 \\ #2 \end{pmatrix}}
\newcommand{\smallColVec}[2]{\left(\begin{smallmatrix} #1 \\ #2 \end{smallmatrix}\right)}
\newcommand{\IdMat}{1}
\newcommand{\transpose}[1]{{#1}^{\mathsf{T}}}
\newcommand{\invtranspose}[1]{{#1}^{-\mathsf{T}}}
\newcommand{\adjugate}[1]{#1^{\mathrm{adj}}}
\newcommand{\DetPair}[2]{\det(#1\mid #2)}
\newcommand{\Sform}{S}
\newcommand{\Diagonal}[2]{\mathrm{diag}(#1,#2)}
\newcommand{\MobiusAction}[2]{#1\cdot #2}

% ---------------------------------------------------------------------
% Ambient groups
% ---------------------------------------------------------------------
% G^+ = GL_2(A_f)/Q_{>0}: the primary ambient group.
\newcommand{\AdelicGroup}{\mathcal{G}^{+}}
% G = GL_2(A_f)/Q^*: the group acting faithfully on the modular tower.
\newcommand{\AdelicGroupCoarse}{\mathcal{G}}
% GL_2(Q)^+/Q_{>0}: the rational lift.
\newcommand{\RationalLiftGroup}{\GLplus{2}{\QQ}/\PositiveRationals}
% PGL_2(Q)^+: the classical (projective) group.
\newcommand{\RationalProjGroup}{\PGLplus{2}{\QQ}}
\newcommand{\AdelicProjection}{\pi}
\newcommand{\CoarseningMap}{\varpi}
\newcommand{\ProjectionToPGL}{\rho}

% Versions of a group: adelic (in G^+), lifted (in GL_2(Q)^+/Q_{>0}),
% plain (in PGL_2(Q)^+), and SL_2 version.
\newcommand{\adelic}[1]{\widehat{#1}}
\newcommand{\lifted}[1]{\widetilde{#1}}
\newcommand{\slversion}[1]{\lifted{#1}\cap\SL{2}{\QQ}}

% Decoration switch for the standard groups: [ad], [lift] or [] (plain).
\newcommand{\GammaSymbol}{\Gamma}
\newcommand{\grpdecor}[2]{\ifstrequal{#1}{ad}{\adelic{#2}}{\ifstrequal{#1}{lift}{\lifted{#2}}{#2}}}
\newcommand{\GammaZero}[2][]{\grpdecor{#1}{\GammaSymbol}_{0}(#2)}
\newcommand{\GammaUpperZero}[2][]{\grpdecor{#1}{\GammaSymbol}^{0}(#2)}
\newcommand{\GammaOne}[2][]{\grpdecor{#1}{\GammaSymbol}_{1}(#2)}
\newcommand{\GammaFull}[2][]{\grpdecor{#1}{\GammaSymbol}(#2)}
\newcommand{\GammaZeroStar}[2][]{\grpdecor{#1}{\GammaSymbol}_{0}(#2)^{*}}
\newcommand{\GammaZeroPlus}[3][]{\grpdecor{#1}{\GammaSymbol}_{0}(#2)+#3}
\newcommand{\GammaConway}[3][]{\grpdecor{#1}{\GammaSymbol}_{0}(#2\,|\,#3)}
\newcommand{\GammaConwayStar}[3][]{\grpdecor{#1}{\GammaSymbol}_{0}(#2\,|\,#3)^{*}}
\newcommand{\GammaPrincipal}[3][]{\grpdecor{#1}{\GammaSymbol}(#2;#3)}
\newcommand{\GammaOfLevel}[2][]{\grpdecor{#1}{\GammaSymbol}_{#2}}
\newcommand{\BorelMod}[1]{B(#1)}
\newcommand{\CommutatorSubgroup}[1]{[#1,#1]}

% ---------------------------------------------------------------------
% Group theory
% ---------------------------------------------------------------------
\DeclareMathOperator{\Stab}{Stab}
\DeclareMathOperator{\Aut}{Aut}
\DeclareMathOperator{\Hom}{Hom}
\DeclareMathOperator{\Isom}{Isom}
\DeclareMathOperator{\Gal}{Gal}
\DeclareMathOperator{\Ker}{ker}
\DeclareMathOperator{\Img}{im}
\DeclareMathOperator{\Coker}{coker}
\DeclareMathOperator{\Tr}{tr}
\DeclareMathOperator{\CoreOp}{core}
\newcommand{\Normalizer}[2]{N_{#1}(#2)}
\newcommand{\Centralizer}[2]{C_{#1}(#2)}
\newcommand{\Core}[2]{\CoreOp_{#1}(#2)}
\newcommand{\ConjClasses}[1]{\mathrm{Cl}(#1)}
\newcommand{\ConjClassesBy}[2]{\mathrm{Cl}_{#1}(#2)}
\newcommand{\ConjClassOf}[1]{[#1]}
\newcommand{\Index}[2]{[#1:#2]}
\newcommand{\Closure}[1]{\overline{#1}}
\newcommand{\FibreProduct}[3]{#1\times_{#3}#2}
\newcommand{\Fibre}[2]{#1^{(#2)}}
\newcommand{\Diag}[1]{#1_{\mathrm{diag}}}
\newcommand{\CentralizerImage}[1]{S(#1)}
\newcommand{\FibreAction}{\ast}
\newcommand{\ModularAut}[1]{\Aut_{\mathrm{mod}}(#1)}
\newcommand{\GeomModularAut}[1]{A(#1)}

% ---------------------------------------------------------------------
% Valuations, divisors, types
% ---------------------------------------------------------------------
\newcommand{\val}[2]{v_{#1}(#2)}
\newcommand{\NumPrimeFactors}[1]{\omega(#1)}
\newcommand{\PsiIndex}[1]{\psi(#1)}
\newcommand{\EulerPhi}[1]{\varphi(#1)}
\newcommand{\HallDivides}{\,\Vert\,}
\newcommand{\TypeGroup}[1]{T_{#1}}
\newcommand{\TypeGroupAll}{T}
\newcommand{\TypeMap}{\mathrm{t}}
\newcommand{\TypeProduct}[2]{#1\ast #2}
\newcommand{\ALimage}[1]{W_{#1}}
\newcommand{\ALdim}[1]{d(#1)}
\newcommand{\DetClass}{\delta}

% ---------------------------------------------------------------------
% The big picture
% ---------------------------------------------------------------------
\newcommand{\Tree}[1]{T_{#1}}
\newcommand{\TreeRealization}[1]{|T_{#1}|}
\newcommand{\BigPicture}{\mathcal{B}}
\newcommand{\BigPictureVertices}{\mathcal{V}}
\newcommand{\LocalVertexSet}[1]{\mathcal{V}_{#1}}
\newcommand{\LatticeClass}[1]{[#1]}
\newcommand{\BaseVertex}{o}
\newcommand{\LocalBaseVertex}[1]{o_{#1}}
\newcommand{\StdLattice}[1]{L_{#1}}
\newcommand{\StdVertex}[1]{o(#1)}
\newcommand{\LocalStdVertex}[2]{o(#1)_{#2}}
\newcommand{\HyperDistance}[2]{\delta(#1,#2)}
\newcommand{\TreeDistance}[3]{d_{#1}(#2,#3)}
\newcommand{\Midpoint}[2]{m(#1,#2)}
\newcommand{\Geodesic}[2]{[#1,#2]}
\newcommand{\FixedSet}[1]{\mathrm{Fix}(#1)}
\newcommand{\LocalFixedSet}[2]{\mathrm{Fix}_{#1}(#2)}
\newcommand{\ConvexHull}[1]{\mathrm{hull}(#1)}
\newcommand{\Centre}[1]{c(#1)}
\newcommand{\LocalCentre}[2]{c_{#1}(#2)}
\newcommand{\StdConfig}[1]{c_{#1}}
\newcommand{\MaxCompact}[1]{\mathcal{U}_{#1}}
\newcommand{\Snake}[1]{\mathcal{S}_{#1}}
\newcommand{\Diameter}[1]{\mathrm{diam}(#1)}

% ---------------------------------------------------------------------
% Invariants of (centred) groups
% ---------------------------------------------------------------------
\newcommand{\HellingLevel}[1]{\operatorname{hl}(#1)}
\newcommand{\CentredLevelQuotient}[1]{\operatorname{m}(#1)}
\newcommand{\CentredLevel}[1]{\operatorname{cl}(#1)}
\newcommand{\QuotientGroup}[2]{Q_{#1,#2}}
\newcommand{\LiftedQuotientGroup}[2]{\widetilde{Q}_{#1,#2}}
\newcommand{\AdelicQuotientGroup}[2]{\widehat{Q}_{#1,#2}}

% ---------------------------------------------------------------------
% Orders, Atkin-Lehner elements, diamond operators
% ---------------------------------------------------------------------
\newcommand{\EichlerOrder}[1]{R_{0}(#1)}
\newcommand{\AdelicEichlerOrder}[1]{\widehat{R}_{0}(#1)}
\newcommand{\IwahoriOrder}[1]{\Lambda_{#1}}
\newcommand{\IwahoriRadical}[1]{J_{#1}}
\newcommand{\Iwahori}[1]{I_{#1}}
\newcommand{\NormalizerIwahori}[1]{\mathcal{N}_{#1}}
\newcommand{\ALmatrix}[1]{W_{#1}}
\newcommand{\ALlocal}[1]{w_{#1}}
\newcommand{\ALflip}[1]{w_{#1}}
\newcommand{\ALconj}[1]{\phi_{#1}}
\newcommand{\ALcocycle}[2]{\kappa(#1,#2)}
\newcommand{\diamondop}[1]{\langle #1\rangle}
\newcommand{\ModelTuple}[4]{(#1,#2,#3,#4)}

% ---------------------------------------------------------------------
% Modular curves and moduli
% ---------------------------------------------------------------------
\newcommand{\ModCurve}[1]{X_{#1}}
\newcommand{\OpenModCurve}[1]{Y_{#1}}
\newcommand{\ComplexPoints}[1]{#1(\CC)}
\newcommand{\PeriodVector}[1]{w_{#1}}
\newcommand{\LatticeOf}[1]{L_{#1}}
\newcommand{\EllCurve}[1]{E_{#1}}
\newcommand{\LatticeBasis}[1]{b_{#1}}
\newcommand{\StdLevel}[2]{\beta_{#1,#2}}
\newcommand{\Torsion}[2]{#1[#2]}
\newcommand{\WeilPairing}[1]{e_{#1}}
\newcommand{\LevelStructure}{\alpha}
\newcommand{\LevelClass}[2]{#1#2}
\newcommand{\ModuliPoint}[2]{(#1,#2)}
\newcommand{\DoubleCosetFinite}[1]{\GL{2}{\ZZ}\backslash(\SlitPlane)\times\GL{2}{\ResidueRing{#1}}}
\newcommand{\ConstantField}[1]{K_{#1}}
\newcommand{\Artin}{\mathrm{art}}
\newcommand{\Reciprocity}[1]{\sigma_{#1}}
\newcommand{\ModAut}[1]{\alpha_{#1}}
\newcommand{\GaloisSemilinearAut}[2]{\Aut_{#1}(#2)}

% ---------------------------------------------------------------------
% Complex abelian varieties without coordinates (appendix)
% ---------------------------------------------------------------------
\newcommand{\GSp}[2]{\mathrm{GSp}_{#1}\optring{#2}}
\newcommand{\Sp}[2]{\mathrm{Sp}_{#1}\optring{#2}}
\newcommand{\SimilitudeFactor}{\nu}
\newcommand{\LieAlg}[1]{\mathrm{Lie}(#1)}
\newcommand{\Differentials}[1]{H^{0}(#1,\Omega^{1}_{#1})}
\newcommand{\Homology}[2]{H_{1}(#1,#2)}
\newcommand{\Cohomology}[2]{H^{1}(#1,#2)}
\newcommand{\HodgeOneZero}[1]{H^{1,0}(#1)}
\newcommand{\ExpMap}[1]{\exp_{#1}}
\newcommand{\PeriodInclusion}[1]{\iota_{#1}}
\newcommand{\Period}[2]{\int_{#1}#2}
\newcommand{\ComplexStructure}[1]{J_{#1}}
\newcommand{\HodgePiece}[3]{\Homology{#1}{\CC}^{#2,#3}}
\newcommand{\HodgeFiltration}[1]{F^{0}\Homology{#1}{\CC}}
\newcommand{\FrameDifferential}[1]{\omega_{#1}}
\newcommand{\DualFrame}[1]{#1^{\vee}}
\newcommand{\PoincareMap}[1]{\mathrm{PD}_{#1}}
\newcommand{\Complexification}[1]{{#1}_{\CC}}
\newcommand{\Realification}[1]{{#1}_{\RR}}
\newcommand{\IntersectionPairing}[2]{#1\cdot #2}
\newcommand{\RiemannForm}[1]{\psi_{#1}}
\newcommand{\PoincareDual}[2]{u_{#1}(#2)}
\newcommand{\Annihilator}[1]{#1^{\perp}}
\newcommand{\Frame}{b}
\newcommand{\PeriodPoint}[1]{\tau_{#1}}
\newcommand{\SiegelPoint}[1]{Z_{#1}}
\newcommand{\SiegelSpace}[1]{\mathfrak{H}_{#1}}
\newcommand{\StdSymplectic}[1]{S_{#1}}
\newcommand{\PeriodBlock}[1]{P_{#1}}
\newcommand{\SiegelLatticeMap}[1]{b_{#1}}
\newcommand{\TorusOf}[1]{A_{#1}}
\newcommand{\NormalizedDifferential}[2]{\omega_{#1,#2}}
\newcommand{\FreeModule}{\mathcal{H}}
\newcommand{\DivisibleGroup}[1]{\mathcal{D}(#1)}
\newcommand{\DivTorsion}[2]{\mathcal{D}(#1)[#2]}
\newcommand{\AbstractTate}[1]{\widehat{T}(\mathcal{D}(#1))}
\newcommand{\pAdicTate}[2]{T_{#1}(\mathcal{D}(#2))}
\newcommand{\DivisionMap}[1]{\mathrm{div}_{#1}}
\newcommand{\TateModule}[1]{\widehat{T}(#1)}
\newcommand{\RationalTateModule}[1]{V_{f}(#1)}
\newcommand{\IdentityBlock}[1]{1_{#1}}
\newcommand{\Covolume}[1]{\mathrm{covol}(#1)}

% ---------------------------------------------------------------------
% Labels
% ---------------------------------------------------------------------
\newcommand{\LabelString}[1]{\texttt{#1}}
